\section{Discussion} \label{sec:discussion}

\subsection{The Informational Causal Gap: Theoretical Limits of Recoverability}
The core theoretical contribution of this work is the rigorous characterization of the \textbf{Informational Causal Gap (ICG)}. In the burgeoning field of digital phenotyping, the ICG represents the fundamental epistemic limit at which the entropy of the missingness process $H(M)$ consumes the mutual information between the latent psychiatric state $S_t$ and the observable behavioral proxy $Y_t$. While traditional missing data literature often treats data loss as a stochastic nuisance to be marginalized out, we posit that in high-stakes clinical psychiatry, the ``absence of signal'' is itself a high-dimensional observation of the latent manifold.

When a patient enters a state of digital disengagement for a duration $\Delta t$, the uncertainty regarding their latent trajectory grows non-linearly. According to the properties of the underlying Neural Stochastic Differential Equation (SDE), the latent state $S_t$ follows an Ito process $dS_t = f(S_t, t)dt + g(S_t, t)dW_t$. In the absence of observations, the variance of the posterior distribution $P(S_{t+\Delta t} | \mathcal{H}_t, M_{t:\Delta t}=1)$ is governed by the diffusion coefficient $g(S_t, t)$. We formally define the \textit{Critical Event Horizon} $\tau^*$ as the temporal threshold beyond which the model's predictive power falls below a clinically actionable safety margin:
\begin{equation}
    \tau^* = \inf \{ \Delta t : \text{Var}(S_{t+\Delta t} | \mathcal{H}_t, M_{t:\Delta t}=1) > \epsilon \}
\end{equation}
where $\mathcal{H}_t$ represents the historical trajectory of both observed behaviors and missingness patterns, and $\epsilon$ is a domain-specific clinical safety threshold (e.g., the maximum allowable variance before a suicidal ideation forecast becomes non-identifiable).

Our empirical results suggest that while smooth diffusion models (e.g., standard Neural SDEs or CSDI) have a $\tau^*$ of approximately 2--3 days in psychiatric contexts, the inclusion of the MNAR-aware Jump-Diffusion prior extends this horizon to 8--10 days. This extension is achieved by leveraging the \textit{temporal velocity} of disengagement as a predictive feature. Specifically, the SDE drift function $f(\cdot)$ is conditioned on the missingness indicator $M_t$, allowing the model to distinguish between a ``random walk'' in the latent space and a directed state transition towards crisis. However, the ICG remains a hard limit. As $\Delta t$ approaches $\tau^*$, the latent state distribution $P(S_{t+\Delta t} | \dots)$ approaches the stationary distribution of the population prior, effectively erasing the patient's individual ``N-of-1'' signal. We argue that the primary role of AI in safety-critical psychiatry should not be to provide a ``best guess'' indefinitely, but to quantify this gap in real-time and signal the exact moment when the model is no longer medically competent to monitor the patient.

\subsection{The Digital Vacuum Paradox: Why Silence is a Primary Signal}
The field of digital psychiatry is currently grappling with what we term the \textbf{``Digital Vacuum Paradox''}. Traditional machine learning assumes that data richness is the prerequisite for insight. In psychiatry, however, the most critical clinical events—such as acute depressive withdrawal, psychotic decompensation, or suicide attempts—are frequently characterized by a total cessation of digital interaction. This creates a paradox where the ``Digital Smoke Detector'' \cite{nock2026predicting} is silenced precisely when the fire is most intense.

\subsubsection{The Myth of Stationarity and the Failure of Traditional Imputation}
Traditional imputation methods (e.g., GAIN, BRITS, LOCF) operate under a subtle assumption of temporal stationarity or slow-moving dynamics. They treat missing values as essentially hidden versions of the observed values, subject to some random or systematic dropout (MAR). They attempt to ``smooth over'' the gaps, effectively interpolating a straight line through a crisis. Our results demonstrate that this smoothing leads to a catastrophic \textit{Lagging Bias}. Because these models do not recognize the disengagement as a signal of a state transition, they only realize a crisis has occurred once the patient re-emerges and generates a new, low-score data point. By then, the intervention window has often closed.

In contrast, our Clinical DLP framework treats $M_t = 1$ (the event of data loss) as a primary feature within the drift function of the SDE. We move from asking ``What would the GPS coordinates have been?'' to ``What does the fact that the GPS is off tell us about the patient's latent risk?'' By modeling the missingness process as a jump process ($dN_t$) that is coupled to the clinical state, we transform the vacuum into a signal. This architectural shift allows the model to ``anticipate'' the crisis during the blackout, providing the clinician with a warning before the patient has even re-engaged with their device.

\subsection{Human-AI Teaming and the 'Learning to Defer' Paradigm}
A central pillar of the Clinical DLP framework is its commitment to the \textbf{Human-AI Teaming} model. We explicitly reject the notion of fully autonomous AI in high-risk psychiatry. Instead, we implement a \textbf{Learning to Defer (L2D)} architecture \cite{joshi2021learning, mozannar2020consistent}, where the system's primary output is not just a risk score, but a decision: ``Predict'' or ``Defer to Human Navigator.''

\subsubsection{Implementing the Clinical Vigilance Index (CVI)}
The engine of this decision is the \textbf{Clinical Vigilance Index (CVI)}. The CVI is a composite metric that quantifies the total risk across two axes: Aleatoric Risk (predicted clinical severity) and Epistemic Risk (uncertainty due to the data gap). Specifically, we define:
\begin{equation}
    CVI_t = \text{logit}^{-1} \left( \alpha \cdot \mathbb{E}[S_t | \mathcal{H}_t] + \beta \cdot \text{KL}(P(S_t | \mathcal{H}_t) || P(S_{t} | \text{pop})) + \gamma \cdot P(M_t | S_{t-1}, X_t) \right)
\end{equation}
where the second term represents the information loss due to disengagement and the third term incorporates socio-technical covariates.

Consider a clinical vignette: A patient with a history of Bipolar Disorder stops transmitting data for 48 hours. A traditional model might report a ``Moderate Risk'' based on the last known stable state. However, the DLP-SDE observes that this disengagement happened at 2:00 AM following a period of high-velocity mobility (potential manic flight). The CVI spikes not just because of the risk score, but because the \textit{Information Gain} regarding the latent state has plummeted. The system defers the case to a human navigator with an alert: ``Informational Causal Gap exceeds safety threshold; high-velocity disengagement detected.'' This allows the navigator to prioritize this silent patient over another who is actively transmitting stable, low-risk data.

\subsection{Socio-Technical Equity: Decoupling Technical Dropout from Clinical Crisis}
A major ethical challenge in digital phenotyping is the risk of \textbf{Digital Exclusion Bias} and what we term \textbf{``Digital Redlining''} \cite{schultebraucks2025llms}. If a model naively interprets every data gap as a sign of clinical disengagement, it will systematically over-identify risk in populations with limited technical resources. Patients from lower socio-economic backgrounds often use older smartphone hardware with degraded batteries, have restricted data plans that lead to frequent connectivity blackouts, or live in areas with poor cellular density.

\subsubsection{Critique of Hardware Bias and Battery Throttling}
Our research identifies several distinct ``Technical Signatures'' that must be decoupled from clinical signals to ensure equity:
\begin{itemize}
    \item \textbf{Battery Throttling Signature:} Modern operating systems aggressively terminate background sensor processes when battery levels fall below a critical threshold (e.g., \SI{15}{\percent}). We observed that trajectories from older devices (e.g., iPhone 8, budget Android models) exhibit ``choppy'' data streams that look remarkably similar to the lethargic behavior associated with MDD. Without explicit hardware calibration, the model misinterprets the device's struggle to stay alive as the patient's struggle to stay active.
    \item \textbf{Connectivity Jitter and Shredded Trajectories:} Frequent switching between 3G/LTE and weak WiFi in rural areas introduces high-frequency noise. Our SDE handles this by treating jitter as a high-frequency noise term $\sigma dW_t$ rather than a structural jump $dN_t$, preventing false positives for behavioral instability.
    \item \textbf{Storage-Induced Silence:} Low-end devices frequently fail to cache sensor data if a large background task (like a WhatsApp backup) is running. This creates periodic ``black holes'' that the DLP-SDE correctly identifies as technical failures by observing the device's telemetry $X_t$.
\end{itemize}

By conditioning the missingness probability on battery health, signal strength, and hardware age, we are able to ``subtract'' the technical noise. We observed that without this calibration, the CVI was \num{1.4}x more likely to flag a crisis in simulated trajectories with ``High Hardware Noise.'' True clinical equity in the digital age requires a framework that is as sensitive to the hardware as it is to the human, ensuring that vigilance is applied equitably across the digital divide.

\subsection{Causal Transportability and N-of-1 Personalization}
While our current benchmarking utilized a large-scale synthetic cohort, the ultimate application of the DLP framework is in the \textbf{N-of-1} setting. Each patient has a unique ``missingness fingerprint.'' For some, disengagement is a sign of a looming crisis; for others, it is a sign of healthy recovery and a return to real-world social engagement (e.g., leaving the phone behind to go for a walk).

The use of Neural SDEs allows us to personalize the model's priors over time. As the system observes a patient's historical response to data gaps, it can adapt the jump intensity $\lambda$ and the diffusion coefficient $\sigma$. This process, which we call \textbf{Latent Bayesian Personalization}, allows the model to learn the patient's individual ``baseline silence.'' This moves digital psychiatry beyond the ``one-size-fits-all'' population models and toward a truly individualized causal understanding of mental health.

Furthermore, we must consider the \textbf{Causal Transportability} of these models \cite{bareinboim2016causal}. Dynamics differ drastically between psychiatric phenotypes. In MDD, disengagement is often a ``slow burn'' decay. In Bipolar Disorder (BD), it can be an ``explosive'' jump. Our framework's ability to condition the SDE drift on diagnostic priors ensures that the recovery of latent states remains clinically valid across diverse phenotypes.

\subsection{Future Directions: Multi-Sensor Fusion and Privacy Shields}
The next generation of the Clinical DLP framework will focus on scaling from single-sensor monitoring to \textbf{Multi-Sensor Fusion}. We envision the integration of physiological signals—such as heart rate variability (HRV) from wearables or even ocular biomarkers (Oculomics). The inclusion of foundation models like \textbf{RETFound} \cite{zhou2023retfound} could provide a holistic ``Latent Patient Representation'' that bridges the Informational Causal Gap even during total smartphone blackouts.

Additionally, we are exploring the use of **Synthetic Cohorts as Privacy Shields**. By training our DLP models on differentially private synthetic data (DSCs), we can share robust recovery architectures across institutions without ever exposing raw, sensitive patient trajectories. This ``Federated Resilience'' would allow a clinic in a rural area to benefit from the complex missingness priors learned in a high-density urban hospital.

\subsection{Ethical Governance and Synthetic Data Integrity}
The use of generative models (Neural SDEs) to reconstruct a patient's ``missing'' life raises significant ethical and legal questions. We must be transparent about the fact that a ``recovered'' trajectory is a statistical estimate, not a recorded fact. We advocate for a \textbf{Transparency First} approach:
\begin{enumerate}
    \item \textbf{Explicit Labeling:} Any data point reconstructed via the SDE must be permanently tagged in the medical record.
    \item \textbf{Patient Sovereignty:} Patients should have the right to review and contest their ``recovered'' data, ensuring the AI-generated version aligns with their lived experience.
    \item \textbf{Algorithmic Auditability:} Regular audits of the missingness model to ensure decoupling of technical and clinical dropout remains effective across demographic groups.
\end{enumerate}

\subsection{Limitations and Adversarial Vulnerabilities}
Despite the robustness of the Jump-Diffusion prior, several limitations remain. First, the framework assumes a degree of temporal consistency. If a patient experiences a \textbf{``Black Swan'' Event}—a total behavioral shift entirely uncorrelated with their history (e.g., a sudden external trauma)—the SDE will fail to recover it. Second, we must consider \textbf{Adversarial Disengagement}. As patients become aware of monitoring, they may intentionally manipulate engagement patterns (e.g., leaving the phone home while engaging in high-risk behavior) to avoid detection. Future research should explore game-theoretic models to ensure the DLP framework remains robust to such strategic behavior.

\section{Conclusion} \label{sec:conclusion}
The evolution of digital psychiatry from a reactive to a proactive discipline depends entirely on our ability to maintain visibility during a patient's most difficult moments. The ``Digital Vacuum Paradox'' has long been the Achilles' heel of the field, where the very signals we seek to measure are the first to vanish in a crisis.

This work has presented a novel \textbf{Clinical Data Loss Prevention (DLP) Framework} that moves beyond the treatment of missing data as a nuisance. By integrating Neural Stochastic Differential Equations with a causally-grounded missingness prior, we have shown that it is possible to transform digital silence into an informative clinical signal. Our framework not only recovers latent psychiatric transitions with \SI{22}{\percent} greater accuracy than state-of-the-art baselines but also introduces the \textbf{Clinical Vigilance Index (CVI)}—a critical tool for navigating the uncertainty of the informational causal gap.

As we move toward a future of safe, resilient, and equitable digital phenotyping, the principles of human-AI teaming, socio-technical calibration, and N-of-1 personalization must remain at the forefront. The Clinical DLP framework ensures that when a patient's digital heartbeats stop, the clinical gaze does not. By quantifying what is knowable and flagging what is not, we provide a foundation for a new generation of psychiatric care that is both technically sophisticated and deeply human.
